\documentclass[
    aps,
    prx,
    reprint,
    superscriptaddress,
    amsmath,
    amssymb,
    longbibliography
]{revtex4-2}

\usepackage{graphicx}
\usepackage{bm}
\usepackage{hyperref}
\usepackage{microtype}
\usepackage{booktabs}
\usepackage{enumitem}

\usepackage{algorithm}
\usepackage{algpseudocode}

\begin{document}

\title{Notes on Reinforcement Learning for Quantum Sensing: Phase and Field Estimation Across Quantum Platforms}

\author{Morten Hjorth-Jensen }
\affiliation{QTAI and UiO, Norway}

\date{\today}

\begin{abstract}
Reinforcement learning (RL) has emerged as a powerful paradigm for the
optimization of quantum sensors operating in complex, noisy, and
time-dependent environments. This article reviews the application of RL
to quantum sensing tasks with an emphasis on phase and field estimation.
We focus on platforms relevant to contemporary quantum technology,
including quantum dots, trapped ions, and neutral atoms. The theoretical
framework combines Bayesian inference, Fisher information, and
partially observable Markov decision processes (POMDPs), allowing RL
agents to adaptively design probe states, control sequences, and
measurement strategies. We discuss model-free and model-aware RL
approaches, provide pseudocode for representative algorithms, and
analyze practical considerations for experimental integration. The
results demonstrate how RL can autonomously approach quantum-limited
sensitivity and offers a unifying methodology for next-generation
quantum metrology.
\end{abstract}

\maketitle

\section{Introduction}

Quantum sensing exploits the coherent dynamics of quantum systems to
estimate physical parameters such as phases, frequencies, and fields
with extreme precision. In idealized scenarios, the achievable
precision is bounded by the quantum Cram\'er--Rao inequality, which is
determined by the quantum Fisher information (QFI) of the probe state.
In practice, however, realistic sensors must operate under noise,
decoherence, control imperfections, and incomplete information about
the parameter of interest.

A canonical sensing task is phase estimation: preparing a quantum
state, letting it pick up an unknown phase $\phi$, then measuring to
infer $\phi$.  The ultimate precision is limited by the Cramér–Rao
bound, $\mathrm{Var}(\hat\phi)\ge1/(N F)$, where $F$ is the Fisher
information and $N$ the number of trials.  In quantum metrology, one
maximizes the quantum Fisher information (QFI) of the probe state to
approach the quantum Cramér–Rao limit.  For a pure probe state
$\ket{\psi(\phi)}$
\[
F_Q(\phi)
= 4\Bigl(\langle \partial_\phi \psi|\partial_\phi \psi\rangle - |\langle \psi|\partial_\phi \psi\rangle|^2\Bigr),
\]
which sets a lower bound $\mathrm{Var}(\hat\phi)\ge 1/[N F_Q(\phi)]$.
In practice, one performs adaptive measurements and Bayesian updating
to approach this bound  .  For example, given a prior distribution
$P(\phi)$ and a measurement outcome $m$ with likelihood $P(m|\phi)$,
Bayesian inference yields the posterior $P(\phi|m)\propto
P(m|\phi),P(\phi)$.  Classical Fisher information for a measurement
${p_k(\phi)}$ is

\[
F_C(\phi)=\sum_k \frac{1}{p_k(\phi)}\Bigl(\frac{\partial p_k(\phi)}{\partial \phi}\Bigr)^2,
\]
while the QFI is the maximal $F_C$ over all measurements .

Traditional quantum control methods (GRAPE, CRAB) can optimize fixed
sensing tasks, but become expensive or fragile for complex, noisy,
time-dependent scenarios . Machine learning offers an alternative: RL
can train an agent to design control protocols or adaptive measurement
policies through trial-and-error.  Recent works have demonstrated deep
RL (policy gradients, Q-learning, cross-entropy) for quantum control
and sensing, often exceeding conventional strategies.  For example,
RL-optimized control pulses in a “quantum-chaotic” sensor were shown
to reduce decoherence effects and improve sensitivity by an order of
magnitude.  Model-aware RL (using gradient information via automatic
differentiation) has also been applied, achieving state-of-the-art
Bayesian experimental design across platforms. This report expands on
these approaches for specific sensor platforms and provides practical
guidance for experimenters.



Adaptive sensing strategies based on Bayesian inference have proven
effective in approaching quantum limits, but the design of optimal
control and measurement policies quickly becomes intractable as the
complexity of the sensor increases. Reinforcement learning provides a
natural framework for addressing this challenge by enabling an agent to
learn adaptive strategies through interaction with the sensing
environment.

This work surveys RL-based approaches to quantum sensing, emphasizing
phase and field estimation and highlighting concrete implementations in
quantum dots, trapped ions, and neutral atom systems.



\section{Bayesian Updates and POMDP Formulation}

Adaptive quantum sensing can be formulated as a partially observable
Markov decision process (POMDP). The unknown parameter $\theta$ is treated
as a hidden state, while the agent maintains a belief distribution
$b(\theta)$ updated via Bayes' rule.
The agent’s belief state is the
Bayesian posterior $b(\theta)$ over $\theta$.  Actions $a$ correspond
to experimental controls (choice of probe state, evolution time,
control Hamiltonians, or measurement basis), and observations $o$ are
measurement outcomes.  The transition dynamics are trivial (the
parameter does not change between measurements), but each action $a$
yields an outcome $o$ with probability $P(o|\theta,a)$. After
observing $o$, the belief is updated via Bayes’ rule:
\begin{equation}
b'(\theta)=
\frac{P(o|\theta,a)\,b(\theta)}
{\int d\theta'\,P(o|\theta',a)\,b(\theta')},
\end{equation}
where $a$ denotes a control action and $o$ a measurement outcome.

The reward function is typically chosen to reflect information gain or
precision, such as the reduction in posterior variance or entropy.
Representing the belief state explicitly converts the POMDP into a fully
observable belief-space MDP, enabling the use of standard RL algorithms.


\section{Fisher Information and Quantum Limits}

For a measurement distribution $p(o|\theta)$, the classical Fisher
information is
\begin{equation}
F_C(\theta)=\sum_o \frac{1}{p(o|\theta)}
\left(\frac{\partial p(o|\theta)}{\partial \theta}\right)^2 .
\end{equation}
The quantum Fisher information $F_Q$ provides an upper bound on $F_C$
over all possible measurements. For pure states,
\begin{equation}
F_Q = 4\left(
\langle \partial_\theta \psi | \partial_\theta \psi \rangle
- |\langle \psi | \partial_\theta \psi \rangle|^2
\right).
\end{equation}

RL policies can be trained to maximize rewards related to Fisher
information, thereby driving the sensing protocol toward
quantum-limited performance.

\section{Reinforcement Learning Methods}





\subsection{Model-Free Reinforcement Learning}

Model-free RL methods (e.g. Q-learning, Deep Q-Networks, Policy
Gradient/A3C, Proximal Policy Optimization) treat the sensor as a
black-box. The agent learns via trial outcomes and rewards, without
explicit knowledge of the system model. Such methods can handle highly
nonlinear, noisy dynamics and accommodate arbitrary reward functions 
. However, they may require large numbers of episodes to converge,
which can be a limitation in experimental settings. For example,
earlier RL approaches to phase estimation achieved good results but
needed many iterations .


\subsection{Model-Aware Reinforcement Learning}

Model-aware RL incorporates known physics into the learning. If a
(differentiable) model of the sensor dynamics is available, one can
backpropagate gradients through the simulation to optimize a policy
network. Belliardo et al. demonstrated that such gradient-based RL
combined with particle-filter Bayesian updates can automatically
design near-optimal experimental sequences across platforms 
. Model-aware methods can directly optimize a reward defined in terms
of Fisher information or estimation variance . They also allow
leveraging analytical gradients to speed up training.


\subsection{Offline RL and Domain Randomization}

Offline RL and Domain Randomization: In many quantum experiments, data
are expensive or time-consuming to acquire. Offline RL (learning from
a fixed dataset) can exploit pre-collected measurements or simulated
data to train policies without active experimentation. Domain
randomization extends this by intentionally varying simulation
parameters (noise levels, decoherence rates, system inhomogeneities)
during training, so that the learned policy is robust to model
mismatch  . For instance, one might randomize a simulated magnetic
field or phase shift to ensure the agent performs well under
calibration uncertainties. While specific work on offline RL in
quantum sensing is still emerging, these techniques from robotics and
control could be adapted to quantum setups.


\section{RL Implementation Blueprint}
An actual implementation would follow these steps:
\begin{enumerate}
  \item {\bf Define the environment:} Choose the sensor platform and identify its controllable parameters and observables.  Build a simulation of the sensor dynamics (e.g.\ master equation or Monte Carlo) that includes the unknown parameter. If possible, validate the simulation against experiments.
  \item {\bf State representation:} Decide what information to feed to the agent.  This may be raw data (fluorescence images, spin readouts) or pre-processed features (means and variances). Include memory of previous steps if needed (e.g.\ via recurrent networks) to handle non-Markovian effects.
  \item {\bf Action space:} Specify the set of actions.  For continuous controls, discretization or parameterized actions can be used.  Ensure actions respect experimental limits (e.g.\ maximum pulse amplitude). 
  \item {\bf Reward function:} Choose a reward that captures precision.  A terminal reward $R = -\log(\text{estimation variance})$ is one option.  Alternatively, use stepwise rewards like Bayesian information gain.  Regularize the reward if necessary to avoid trivial strategies.
  \item {\bf RL algorithm:} Select an algorithm based on action type.  For example, use DQN for coarse adjustment of discrete parameters, or SAC/DDPG for smooth continuous control.  Set neural network architectures appropriate for the state (e.g.\ convolutional nets for images).
  \item {\bf Training:} Train the agent in simulation.  Introduce variability (randomize unknowns) to encourage generalization.  Use experience replay or batch updates to improve sample efficiency.  Monitor performance on validation scenarios.
  \item {\bf Deployment:} Transfer the trained policy to the real device.  If training was in-silico, consider fine-tuning on the hardware with a few adjustment episodes.  Include safeguards for safety (e.g.\ bounding pulses) and mechanisms to reset the environment.
\end{enumerate}



\section{Quantum Sensing Platforms}

\subsection{Quantum Dot Sensors}

Semiconductor quantum dot devices can act as tunable charge or spin
sensors for electric and magnetic fields. Double quantum dots (DQDs)
are widely used for qubit readout and have complex stability
diagrams. Deep RL has been applied to automate DQD measurement: Nguyen
et al. demonstrated a DQN-based agent that rapidly locates “bias
triangles” in a transport map of a GaAs double dot . This amounts to
learning an experimental protocol (gate voltages to apply and where to
measure) to identify a narrow feature in a large parameter
space. Their algorithm found triangles in minutes, far faster than
manual tuning . This approach shows that RL can replace expert
heuristics in navigating dot operation regimes. Similar RL agents
could be trained to adaptively measure phase shifts in a quantum dot
interferometer or to optimize gate voltages for maximum charge
sensitivity.

In practice, one would model the dot sensor’s output $o$ (current or
charge) as a function of the unknown parameter $\phi$ (e.g. a phase
imposed on the electron spin) and the control voltages $a$. The agent
then updates its belief on $\phi$ after each measurement. Since dot
charge sensors often have limited quantum coherence, these tasks can
often be well described as classical POMDPs, though spin qubits in
dots could require quantum updates. The RL policy might be a neural
network taking the current belief and possibly past history as input
and outputting the next gate voltage configuration. Domain
randomization can account for device variability between cooldowns.
In short, RL offers a flexible framework for autonomous tuning and
adaptive estimation in quantum dot sensors .


\subsection{Trapped Ion Sensors}

Ion traps confine individual atomic ions (e.g. Yb+, Ca+) using
electric and/or magnetic fields. These platforms are used for quantum
clocks, electric-field sensing, and magnetometry. The control
parameters include laser frequencies/pulses and voltages on segmented
electrodes. The challenges include compensating stray fields and
designing pulse sequences under decoherence. Recent studies have
applied ML and RL in this context. For example, Schmitz et al. used
adaptive algorithms to optimize stray electric field compensation for
a ^171Yb⁺ ion . They compared gradient descent (ADAM) with a
machine-learning optimizer (MLOOP) to maximize ion
fluorescence. Starting from a nearly optimized manual setting, ADAM
improved fluorescence by 78$\%$ and MLOOP achieved 96$\%$ .  This shows
ML’s power in high-dimensional control (44 DC electrodes) to maintain
optimal sensor operation.

In this setup, an RL agent could treat the set of compensation
voltages and/or laser beam parameters as the action space. For
parameter estimation (e.g. sensing a phase imprinted on the ion’s
internal state), the agent would iteratively apply a control (like a
rotation angle) and measure fluorescence, updating a phase
posterior. Cross-entropy RL and policy-gradient methods have been used
to find optimal pulse sequences in trapped-ion sensing. Schuff et
al. demonstrated that RL (cross-entropy method) can discover
spin-squeezing–like pulse strategies in a collective-spin ion sensor,
outperforming simple periodic pulses . Importantly, the agent learned
to mitigate decoherence by adjusting pulse timing and strength,
yielding up to an order-of-magnitude sensitivity gain . This adaptive
RL protocol effectively enhances the quantum Fisher information of the
sensing scheme.

In implementation, one must consider experimental constraints:
real-time feedback bandwidth, laser noise, and ion heating. Offline RL
or simulation-based pretraining can reduce the number of costly live
trials. For instance, by simulating the ion’s master-equation
dynamics, an agent can learn a policy before deployment, then
fine-tune on the actual hardware. Domain randomization (varying
decoherence rates or trap frequencies in simulation) can make the
policy robust to model mismatch. Integration with ion-trap software
(e.g. ARTIQ or custom FPGA control) allows the agent to select pulses
or compensation voltages on the fly. Careful reward shaping
(e.g. using reductions in state variance or improvements in
fluorescence) will guide the RL to physically meaningful sensing
strategies  .


\subsection{Neutral Atom Sensors}

Neutral-atom platforms include cold-atom interferometers and
magnetometers. Alkali vapor cells (e.g. Rb, Cs) and ultracold
Bose–Einstein condensates (BECs) can serve as precision sensors for
fields and acceleration. A recent cutting-edge example is using
ultracold atoms in an optical lattice as an inertial sensor. Chih and
Holland applied reinforcement learning to design a novel atom-lattice
gyroscope  . Instead of the standard Mach–Zehnder sequence, the RL
agent learned an end-to-end “shaking” protocol for a 2D
lattice. Remarkably, for the same total interrogation time, the
learned protocol achieved a 20-fold improvement in rotation
sensitivity over conventional interferometry . This demonstrates RL’s
ability to discover non-intuitive quantum sensing schemes by
maximizing long-term rewards.

Neutral atoms also include solid-state ensembles like NV-center
magnetometers, which can be treated similarly. An agent could
adaptively choose microwave pulses and measurement bases to track a
time-varying magnetic field, updating a field estimate after each
measurement. Bayesian RL approaches (policy gradient with belief
updates) could greatly speed up such magnetometry.

On Earth, atomic ensembles in magneto-optical traps or lattice arrays
act as vector magnetometers and clocks. Reinforcement learning can
optimize these too: for example, an RL agent might learn how long to
interrogate a Ramsey sequence based on the current estimate of an
oscillating field, effectively learning adaptive control that
approaches the quantum projection noise limit. Experiments have begun
to implement RL for multi-atom systems; one may adapt existing RL
toolkits (e.g. qsensoropt ) to these platforms. Key practical
considerations include dealing with large particle numbers and
decoherence. Fortunately, the notion of QFI extends to many-body
states, and RL can be used to approach the best possible scaling
(shot-noise or beyond). Control actions may include global phase
shifts (via magnetic fields) and collective readout
choices. Reinforcement learning can thus serve as a unifying framework
for many neutral-atom sensing modalities  .



\section{Pseudocode Examples}

\begin{algorithm}[t]
\caption{RL-Based Adaptive Phase Estimation}
\begin{algorithmic}[1]
\State Initialize belief $b(\phi)$, policy $\pi$
\For{each episode}
    \For{$t=1$ to $T$}
        \State Select action $a_t=\pi(b(\phi))$
        \State Perform measurement, observe $o_t$
        \State Update belief $b(\phi)$ via Bayes' rule
        \State Compute reward $r_t$
        \State Update policy $\pi$
    \EndFor
\EndFor
\end{algorithmic}
\end{algorithm}


To illustrate how one might implement an RL-based adaptive estimator, we give pseudocode for two representative tasks.

\begin{algorithmic}[1]
\State \textbf{Algorithm: RL-based Adaptive Phase Estimation}
\State {\bf Initialize} prior distribution $P(\phi)$ uniform over $[0,2\pi)$, policy network (or Q-table) $\pi$, and replay buffer (if applicable).
\For{episode = 1 to $N_{\rm ep}$} 
  \State Reset belief $b(\phi)\leftarrow P(\phi)$.
  \For{step = 1 to $T$}
    \State {\bf Action:} choose control parameter $a_t = \pi(b(\phi))$ (e.g.\ phase shift or pulse duration).
    \State {\bf Probe:} apply $a_t$ to sensor, obtain measurement outcome $m_t\in\{0,1\}$.
    \State {\bf Belief update:} use Bayes: $b(\phi)\propto P(m_t|\phi,a_t)\,b(\phi)$, then normalize.
    \State {\bf Reward:} compute reward $r_t$, e.g.\ $-\mathrm{Var}_b(\phi)$ or information gain $D_{\rm KL}(b\parallel b_{\rm prior})$.
    \State Store $(\text{state}=b(\phi),\,\text{action}=a_t,\,r=r_t,\,\text{next state}=b'(\phi))$.
    \State Update policy $\pi$ (Q-learning or gradient step) using the transition and reward.
  \EndFor
\EndFor
\end{algorithmic}

The above loop exemplifies adaptive Bayesian estimation: after each
measurement, the agent updates the posterior $b(\phi)$ and uses it as
input for the next action.  This fully Bayesian RL approach is
data-efficient, reusing all prior information . By shaping the reward
with estimation variance, the agent learns to drive the posterior
toward a narrow distribution, effectively maximizing Fisher
information.

\begin{algorithmic}[1]
\State \textbf{Algorithm: RL for Magnetic Field Estimation with NV Centers}
\State {\bf Initialize} initial state estimate $b(B)$ (over magnetic field $B$), policy $\pi$, and environment model (optional).
\For{episode = 1 to $N_{\rm ep}$}
  \State Reset belief $b(B)$ to prior (e.g.\ broad Gaussian).
  \For{step = 1 to $T$}
    \State {\bf Action:} select control sequence $a_t$ (e.g.\ microwave pulse amplitude and duration, or Ramsey time).
    \State {\bf Apply:} implement $a_t$ on the NV center, record photoluminescence signal $m_t$.
    \State {\bf Likelihood:} compute $P(m_t|B,a_t)$ from the sensor model.
    \State {\bf Belief update:} $b(B)\propto P(m_t|B,a_t)\,b(B)$.
    \State {\bf Reward:} $r_t = -\mathrm{Var}_b(B)$ (or other metric).
    \State Store transition and update policy $\pi$.
  \EndFor
\EndFor
\end{algorithmic}

In both examples, the state of the MDP is the current belief
distribution (or a finite representation of it) and the observation is
the latest measurement outcome.  These pseudocode sketches can be
implemented with deep neural networks (taking a discretized belief or
sufficient statistics as input) or simpler methods if the parameter
space is small.  They illustrate how RL integrates control choice,
Bayesian inference, and reward-driven learning into a closed loop.

\section{Model-Free vs Model-Aware, Offline RL, Domain Randomization}

We have noted that model-free RL (e.g. Q-learning) makes minimal
assumptions but can be data-hungry.  In contrast, model-aware RL (like
policy gradient with differentiable simulation) exploits knowledge of
the sensor’s quantum dynamics.  For instance, Belliardo et
al. designed a library (qsensoropt) that uses automatic
differentiation through the Lindblad evolution and particle filter to
train policies .  Their model-aware agent directly computes gradients
of precision metrics (Fisher information or posterior covariance) with
respect to control parameters, yielding highly efficient training.

Offline RL can leverage pre-recorded experimental data.  For example,
one might perform a set of calibration measurements offline and train
an RL agent on this fixed dataset.  Care must be taken to avoid
overfitting; conservative algorithms (like Batch Constrained
Q-learning) could be adapted. Domain randomization has been crucial in
transferring policies from simulation to real robots; similarly, in
quantum sensing one can randomize Hamiltonian parameters or noise
levels during training.  This way, an RL policy learns to work under
uncertainty (e.g. unknown decoherence rate). While specific
demonstrations in quantum sensing are limited, these methods are
straightforward to implement: one simply varies the simulated
environment dynamics at the start of each training episode. Doing so
can improve robustness when the actual device deviates from the ideal
model.

Finally, hybrid methods like Fiderer et al.’s show that evolutionary
strategies and RL can jointly optimize adaptive heuristics . In some
cases, one could pretrain a neural network to propose near-optimal
actions (via supervised learning on simulated “expert” data) and then
fine-tune it with RL for final performance. This might mitigate the
problem of sparse rewards in experimental data.


\section{Implementation Considerations}

Key practical challenges include data efficiency, reward design,
hardware constraints, and robustness to noise and drift. Hybrid
strategies combining simulation pretraining and experimental fine-tuning
are likely to be essential for scalable deployment.

Integrating RL with real quantum sensors requires attention to practical constraints:
\begin{itemize}
\item    Data Efficiency: Experiments are slow. Model-aware training in simulation is recommended, with periodic calibration on real hardware. Agents can use Gaussian process models or Bayesian surrogate models to reduce sample count.
\item Reward Design: The reward function must reflect sensing goals. Options include negative posterior variance, log-likelihood, or information gain. Explicitly using quantum-limits (e.g. QFI) in the reward can accelerate learning  .
\item Noise and Imperfections: Real measurements have technical noise, drift, and detection errors. Reward functions should penalize noisy actions. Domain randomization during training can help the agent tolerate such imperfections. One can also include a “no-update” action to handle erasure or no-fluorescence outcomes.
\item Hardware Constraints: Action space should respect hardware limits (e.g. discrete pulse durations, limited voltage range). Use realistic simulations of the control hardware. An RL policy network or Q-table can run on real-time controllers (FPGA, GPU) to decide the next move between experimental shots.
\item Exploration vs Exploitation: In an experimental loop, aggressive exploration can degrade performance. A safe strategy is to start with a known good protocol and let RL refine it (e.g. epsilon-greedy around a baseline). Off-policy learning from off-line data can also reduce on-the-fly exploration.
\item Parallel Experiments: If multiple sensors or multiple runs are possible, parallel data collection can speed up training. Care must be taken if agents share a common policy to avoid correlated failures.
\end{itemize}

As a concrete example, consider calibrating an atomic magnetometer: an
RL agent could sample at different interrogation times to build a
field estimate. The pseudocode above shows how each measurement
outcome would refine the field belief. One practical recipe is: (1)
train the RL policy in a high-fidelity simulator that includes
estimated decoherence; (2) load the policy into the experiment and run
in “closed loop” with Bayesian updates; (3) continue occasional
training online using real outcomes. This kind of mixed-model RL
leverages simulations for learning efficiency but corrects for reality
with real data.



\section{Conclusion}

Reinforcement learning provides a versatile and powerful framework for
quantum sensing. By unifying Bayesian inference, Fisher information, and
adaptive control, RL enables autonomous discovery of high-performance
sensing protocols across a wide range of quantum platforms. As quantum
hardware matures, RL-driven sensing is poised to become a central tool
in quantum metrology.

\bibliography{rl_quantum_sensing}

\end{document}






10. Implementation Considerations for Experiments


11. Conclusion

Reinforcement learning provides a powerful and general framework for optimizing quantum sensors. By framing phase and field estimation as POMDPs, RL agents can autonomously adapt probe settings and measurements to maximize information gain. Across platforms—from quantum dots and trapped ions to ultracold atoms—RL has enabled new adaptive protocols and control strategies   . Model-aware RL methods have shown particular promise, achieving or even surpassing traditional limits by leveraging quantum Fisher information  .  Going forward, the synergy of Bayesian estimation, Fisher-information metrics, and reinforcement learning will likely yield the next generation of high-performance quantum sensors. Physicists and engineers implementing RL in the lab should carefully consider simulation fidelity, reward engineering, and hardware constraints as outlined above.  With thoughtful integration, RL can significantly accelerate calibration and control, pushing quantum metrology closer to its ultimate limits.

References
	•	Xiao et al., “Parameter estimation in quantum sensing based on deep reinforcement learning,” npj Quantum Inf. 8, 2 (2022)  .
	•	Schuff et al., “Improving the dynamics of quantum sensors with reinforcement learning,” New J. Phys. 22, 035001 (2020) .
	•	Belliardo et al., “Model-aware reinforcement learning for high-performance Bayesian experimental design in quantum metrology,” Quantum 8, 1555 (2024)  .
	•	Bharti, “Fisher Information: A crucial tool for NISQ research,” Quantum Perspective 5, 61 (2021)  .
	•	Nguyen et al., “Deep reinforcement learning for efficient measurement of quantum devices,” npj Quantum Inf. 7, 100 (2021) .
	•	Schmitz et al., “Dynamic compensation of stray electric fields in an ion trap using machine learning and adaptive algorithm,” Sci. Rep. 12, 20503 (2022) .
	•	Chih & Holland, “Reinforcement Learning for Rotation Sensing with Ultracold Atoms in an Optical Lattice,” arXiv 2212.14473 (2024)  .
	•	Fiderer et al., “Neural-Network Heuristics for Adaptive Bayesian Quantum Estimation,” PRX Quantum 2, 020303 (2021) .
